\section*{Chapter \ref{ch:pl:observability-and-incident-response} --- Observability and Incident Response}

\begin{solution}{exo:pl:observability-and-incident-response:1}
$0.001 \times 30 \times 23\,400 = 702$ bad seconds; the quiet month spent 115, 16.4\% of the budget, in 18 pauses of more than five seconds.
\end{solution}

\begin{solution}{exo:pl:observability-and-incident-response:2}
Because no process stops: in the stall the processes wait for data that does not come, in the intermittent stalls they wait twenty seconds a
minute, in the slow consumer the risk service is busy. A heartbeat measures the process, not the service the users receive.
\end{solution}

\begin{solution}{exo:pl:observability-and-incident-response:3}
56 seconds: the first five seconds are still good (data at most five seconds old); the hour window must then hold 1.44\% of 3\,600 seconds bad,
51.8, so the page comes at the 52nd bad second, with the five-minute window long since past its own limit.
\end{solution}

\begin{solution}{exo:pl:observability-and-incident-response:4}
Each stall lasts twenty seconds and the data then refreshes, so the age never exceeds twenty seconds. A third of the seconds are bad, which
spends the budget quickly, but no single second crosses thirty.
\end{solution}

\begin{solution}{exo:pl:observability-and-incident-response:5}
It requires the problem to be happening now: without it, an hour-long window keeps firing for up to an hour after an incident has ended, and a
burst early in the window keeps paging after recovery. The short window makes the alert stop soon after the service recovers.
\end{solution}

\begin{solution}{exo:pl:observability-and-incident-response:6}
21 minutes, from the last quote at 12:03 to fresh data at 12:24, measured on the service the risk system receives. A mean time to repair
(Book~14, chapter~15) is a component's average from failure to repair over many failures; this is one incident of one service.
\end{solution}

\begin{solution}{exo:pl:observability-and-incident-response:7}
The five-second threshold doubles to 36 false pages a month and the 30-second one triples to 3; the burn-rate rules stay at one, because the
pauses are still short and spend little budget each.
\end{solution}

\begin{solution}{exo:pl:observability-and-incident-response:8}
The services were healthy and idle: they reported on themselves, not on what they delivered. Health checks must be complemented by SLIs from
the users' side --- here the age of the data --- or an outage upstream looks like a quiet day.
\end{solution}

\begin{solution}{pb:pl:observability-and-incident-response:1}
\begin{enumerate}
\item \omterm{def:pl:observability-and-incident-response:metric}{Metrics}: is something wrong, everywhere, cheaply. Logs: what happened in one service. Traces: where the time went across services for one
unit of work.
\item Fractions of a millisecond in capture, store and positions, then the risk queue: 50~ms normally, 60~s during the slow consumer; risk
computes in 1~ms.
\item Because a service can be alive and useless: the users need fresh data, not live processes.
\item 702 bad seconds a month; quiet days spend 115.
\item An SLO is the firm's own target, set tighter; an SLA is a promise to a customer with penalties.
\item A feed stall from 12:03 for 21 minutes, 20-second stalls every minute from 11:48 for 15 minutes, a slow consumer from 14:00.
\item Five seconds: all three, in 5, 5 and 100~s; thirty: the stall in 30~s and the slow consumer in 600~s, not the intermittent stalls;
sixty: the stall in 60~s, the slow consumer in 1\,200~s, not the intermittent stalls.
\item 18, 1 and 0 for the three thresholds, none for the heartbeat, one for the burn-rate rules.
\item All three, in 56, 191 and 151 seconds.
\item It is the 56-second pause, which spends a real share of the budget and deserves a look.
\item Holds the incident's state and assigns the roles; does not do the technical work.
\item Exposure: is the firm trading on stale data, and should strategies stop --- kill switches and limits before diagnosis.
\item 12:03:00 last quote; 12:03:05, 12:03:30, 12:03:56 and 12:04:00 the four pages; 12:24:00 the feed resumes and the data is fresh.
\item 56 seconds to detect under the burn-rate rule, 21 minutes to restore.
\item A timeline, contributing causes without blame, and changes to the system.
\item \textbf{Named result.} A complete feed stall is detected in 5, 30 and 60 seconds by static thresholds at those ages and in 56 seconds by
the burn-rate rules; a slow consumer in 100, 600 and 1\,200 seconds against 151; intermittent stalls only by the five-second threshold (5
seconds) and the burn-rate rules (191). On thirty quiet days the thresholds send 18, 1 and 0 false pages a month, the burn-rate rules one;
a process heartbeat detects nothing.
\item The burn-rate rules for pages, with a ticket on a burn rate of 1 over three days for slow degradation.
\item An SLI on the age of the data each consumer uses.
\item Data-age SLIs everywhere, burn-rate paging, a kill switch on stale data, and the replay kept as a test of the alerting.
\item That a person must act now, on something the users of the service can feel.
\end{enumerate}
\end{solution}

\begin{solution}{iq:pl:observability-and-incident-response:1}
A process check asks whether a program runs; a service check asks whether its users get what they need, measured from their side (data age,
error rate, latency). Only the second sees an upstream failure or a silent slowdown.
\iqlookfor{User-side indicators.}
\end{solution}

\begin{solution}{iq:pl:observability-and-incident-response:2}
An SLI measures the service (the share of seconds with market data at most five seconds old); the SLO is its target (99.9\% over 30 trading
days); the \omterm{def:pl:observability-and-incident-response:budget}{error budget} is what the target allows (702 bad seconds a month), spent on incidents and risky changes.
\iqlookfor{A concrete indicator, target and budget.}
\end{solution}

\begin{solution}{iq:pl:observability-and-incident-response:3}
Delete or demote every page that did not need a person now; page on SLO burn rates instead of raw thresholds; put a \omterm{def:pl:observability-and-incident-response:oncall}{runbook} behind every page;
review the pages each week.
\iqlookfor{Pages mean action.}
\end{solution}

\begin{solution}{iq:pl:observability-and-incident-response:4}
Trace it: propagate a trace identifier with the message through every service, record a \omterm{def:pl:observability-and-incident-response:trace}{span} per step, and read the tree; a queue wait
shows up as one long \omterm{def:pl:observability-and-incident-response:trace}{span}.
\iqlookfor{\omterm{def:pl:observability-and-incident-response:trace}{Distributed tracing}.}
\end{solution}

\begin{solution}{iq:pl:observability-and-incident-response:5}
User-side SLIs and SLOs per service (data age, completeness, latency); \omterm{def:pl:observability-and-incident-response:metric}{metrics}, structured logs and traces with propagated identifiers;
burn-rate pages and slow-burn tickets; \omterm{def:pl:observability-and-incident-response:oncall}{on-call rotations} with \omterm{def:pl:observability-and-incident-response:oncall}{runbooks} that start with exposure; incident command; \omterm{def:pl:observability-and-incident-response:blameless}{blameless reviews} with
tracked actions; replayed incidents as tests.
\iqlookfor{SLOs, burn rates, incident practice.}
\end{solution}
